% !TEX root = ../main.tex
\corrigesChapitre{7}{Vecteurs et produit scalaire}
\corrigesSouples


\begin{blocLivre}{cor-7-1}
\begin{corrige}{1}{Coordonnées et norme}
\demarche{Soustraire les coordonnées dans l’ordre indiqué par le nom du vecteur.}
\begin{questions}
\item $\res{\vect{AB}=(3;4)}$ et $\res{\vect{BA}=(-3;-4)}$.
\item $AB=\sqrt{3^2+4^2}=\res{5}$. Le milieu est $\res{(5/2;4)}$.
\end{questions}

\label{figgeo-07-01-cor}
\figureGeo[corrige]{07-01}
\end{corrige}
\end{blocLivre}
\begin{blocLivre}{cor-7-2}
\begin{corrige}{2}{Un calcul avec les coordonnées}
\demarche{Le produit scalaire utilise les coordonnées ; la norme utilise la somme de leurs carrés.}
\begin{questions}
\item $u\cdot v=6-4=\res{2}$ ; $u\cdot u=5$ ; $v\cdot v=25$.
\item Les normes sont $\res{\sqrt5}$ et $\res{5}$.
\end{questions}

\label{figgeo-07-02-cor}
\figureGeo[corrige]{07-02}
\end{corrige}
\end{blocLivre}
\begin{blocLivre}{cor-7-3}
\begin{corrige}{3}{Longueurs et angle}
\demarche{Le cosinus de l’angle détermine le signe du produit scalaire.}
On utilise $u\cdot v=30\cos\theta$. Les valeurs sont $\res{15}$, $\res{0}$ et $\res{-15}$. Le signe indique si l’angle est aigu, droit ou obtus.
\label{figgeo-07-03-cor}
\figureGeo[corrige]{07-03}
\end{corrige}
\end{blocLivre}
\begin{blocLivre}{cor-7-4}
\begin{corrige}{4}{Orthogonaux ou colinéaires}
\demarche{Tester séparément orthogonalité et colinéarité, qui correspondent à deux critères différents.}
\begin{questions}
\item $u\cdot v=6-6=0$, donc $u\perp v$. $u\cdot w=-18-8=-26\neq0$, donc $u$ et $w$ ne sont pas orthogonaux.
\item $\res{w=-2u}$ : $u$ et $w$ sont colinéaires de sens opposés.
\end{questions}

\label{figgeo-07-04-cor}
\figureGeo[corrige]{07-04}
\end{corrige}
\end{blocLivre}
\begin{blocLivre}{cor-7-5}
\begin{corrige}{5}{Lire une projection}
\demarche{Identifier le projeté sur la droite horizontale puis vérifier le calcul par coordonnées.}
\begin{questions}
\item L’axe $(AB)$ est horizontal : $\res{H(2;0)}$.
\item Par projection, $AB\times AH=6\times2=\res{12}$. Par coordonnées, $(6;0)\cdot(2;3)=12$.
\end{questions}

\label{figgeo-07-05-cor}
\figureGeo[corrige]{07-05}
\end{corrige}
\end{blocLivre}
\begin{blocLivre}{cor-7-6}
\begin{corrige}{6}{Les normes d'une somme}
\demarche{Développer les carrés des normes avec le produit scalaire.}
\begin{questions}
\item Par bilinéarité,
$\|u+v\|^2=\|u\|^2+2u\cdot v+\|v\|^2=9+12+16=\res{37}$.
Avec une différence, le terme croisé change de signe :
$\|u-v\|^2=9-12+16=\res{13}$.
\item Une norme est positive ou nulle : $\res{\|u+v\|=\sqrt{37}}$ et $\res{\|u-v\|=\sqrt{13}}$.
Les carrés des diagonales ont pour somme $37+13=50$, égale à $2\|u\|^2+2\|v\|^2=2\times9+2\times16$.
\end{questions}
\label{figgeo-07-06-cor}
\figureGeo[corrige]{07-06}
\end{corrige}
\end{blocLivre}
\begin{blocLivre}{cor-7-7}
\begin{corrige}{7}{Calculer un angle de triangle}
\demarche{Isoler le cosinus dans la définition du produit scalaire.}
\begin{questions}
\item Le produit vaut $4$, $AB=4$ et $AC=2$.
\item $\cos\widehat{BAC}=4/(4\times2)=1/2$, donc $\res{\widehat{BAC}=\pi/3}$.
\item $BC=\sqrt{(-3)^2+(\sqrt3)^2}=2\sqrt3$. Comme $AC^2+BC^2=4+12=16=AB^2$, le triangle est rectangle en $C$.
\end{questions}

\label{figgeo-07-07-cor}
\figureGeo[corrige]{07-07}
\end{corrige}
\end{blocLivre}
\begin{blocLivre}{cor-7-8}
\begin{corrige}{8}{Déterminer un paramètre}
\demarche{Traduire chaque propriété par une équation sur le paramètre.}
\begin{questions}
\item L’orthogonalité équivaut à $u_m\cdot v=0$.
Or $u_m\cdot v=3m-2$, donc $3m-2=0\iff\res{m=2/3}$.
\item La colinéarité équivaut à un déterminant nul :
$m(-1)-2\times3=-m-6=0\iff\res{m=-6}$.
Dans ce cas $u_{-6}=(-6;2)=-2v$, ce qui contrôle le résultat.
\item Pour $m=0$, le produit vaut $-2$, et les normes valent 2 et $\sqrt{10}$.
La définition donne $\cos\theta=(-2)/(2\sqrt{10})=\res{-1/\sqrt{10}}$.
Ce cosinus négatif correspond bien à un angle obtus.
\end{questions}
\label{figgeo-07-08-cor}
\figureGeo[corrige]{07-08}
\end{corrige}
\end{blocLivre}
\begin{blocLivre}{cor-7-9}
\begin{corrige}{9}{Les diagonales d'un rectangle}
\demarche{Décomposer les diagonales à l’aide des deux côtés du rectangle.}
\begin{questions}
\item $\vect{AC}=u+v$ et $\vect{BD}=v-u$.
\item Comme $u\cdot v=0$, $(u+v)\cdot(v-u)=\|v\|^2-\|u\|^2=\res{-20}$. Puis $u\cdot(u+v)=\res{36}$.
\item Le produit des diagonales n’est pas nul : elles ne sont pas perpendiculaires.
\end{questions}

\label{figgeo-07-09-cor}
\figureGeo[corrige]{07-09}
\end{corrige}
\end{blocLivre}
\begin{blocLivre}{cor-7-10}
\begin{corrige}{10}{Un parallélogramme non rectangle}
\demarche{Calculer d’abord le produit des côtés, puis les carrés des diagonales.}
\begin{questions}
\item Le produit vaut $3\times4\times1/2=\res{6}$.
\item $AC^2=9+16+12=37$ et $BD^2=9+16-12=13$, donc $\res{AC=\sqrt{37},\ BD=\sqrt{13}}$.
\item $37+13=50=2\times9+2\times16$ : l’identité du parallélogramme est vérifiée.
\end{questions}

\label{figgeo-07-10-cor}
\figureGeo[corrige]{07-10}
\end{corrige}
\end{blocLivre}
\begin{blocLivre}{cor-7-11}
\begin{corrige}{11}{Projeter sans coordonnées}
\demarche{Relier le produit scalaire à la projection avant d’utiliser Pythagore.}
\begin{questions}
\item $8\times5\times1/2=\res{20}$.
\item $8\times AH=20$, donc $AH=5/2$. Dans $AHC$, $CH=\sqrt{25-25/4}=\res{5\sqrt3/2}$.
\item $BC^2=64+25-40=49$, donc $\res{BC=7}$.
\end{questions}

\label{figgeo-07-11-cor}
\figureGeo[corrige]{07-11}
\end{corrige}
\end{blocLivre}
\begin{blocLivre}{cor-7-12}
\begin{corrige}{12}{Un triangle et sa hauteur}
\demarche{Vérifier les angles droits par les produits scalaires, puis calculer les longueurs.}
\begin{questions}
\item $\vect{CA}=(-1;-3)$ et $\vect{CB}=(3;-1)$ ont un produit $-3+3=0$.
\item $\vect{AB}=(4;2)$, $\vect{AC}=(1;3)$ : produit $10$, normes $\sqrt{20}$ et $\sqrt{10}$. Le cosinus vaut $1/\sqrt2$, donc l’angle est $\pi/4$.
\item $H=(3;2)$ et $\vect{CH}=(1;-2)$, donc $\vect{CH}\cdot\vect{AB}=4-4=0$.
\item $CH=\sqrt5$. L’aire vaut $AB\times CH/2=\sqrt{20}\sqrt5/2=\res{5}$.
\end{questions}

\label{figgeo-07-12-cor}
\figureGeo[corrige]{07-12}
\end{corrige}
\end{blocLivre}
\begin{blocLivre}{cor-7-13}
\begin{corrige}{13}{Le cercle de diamètre donné}
\demarche{Comparer la distance au centre avec le rayon pour interpréter le signe du produit.}
\begin{questions}
\item $(-2-x)(2-x)+(-y)(-y)=\res{x^2+y^2-4}$.
\item Le produit est nul pour $x^2+y^2=4$, cercle de centre $O$ et de rayon 2.
\item Il est négatif à l’intérieur strict du disque, positif à l’extérieur strict.
\item Pour $P$, il vaut 0 : $P$ est sur le cercle. Pour $Q$, il vaut $-2$ : $Q$ est à l’intérieur.
\end{questions}

\label{figgeo-07-13-cor}
\figureGeo[corrige]{07-13}
\end{corrige}
\end{blocLivre}
\begin{blocLivre}{cor-7-14}
\begin{corrige}{14}{Travail d'une force}
\demarche{Écrire le déplacement avec ses deux composantes et garder les unités physiques.}
\begin{questions}
\item Pour le déplacement horizontal, $W=(30;40)\cdot(6;0)=30\times6+40\times0=\res{180\text{ J}}$.
La norme de la force est $\sqrt{30^2+40^2}=\res{50\text{ N}}$.
\item Un déplacement de norme 6 et d’angle $\pi/6$ a pour coordonnées
$(6\cos(\pi/6);6\sin(\pi/6))=(3\sqrt3;3)$.
Ainsi $W=30(3\sqrt3)+40\times3=\res{90\sqrt3+120\text{ J}}$.
\item Si $\theta$ est l’angle entre la force et le déplacement,
$W=50\times6\cos\theta\leq300$ puisque $\cos\theta\leq1$.
L’égalité est atteinte dans la direction et le sens de $F$ :
$d=6F/\|F\|=6(30;40)/50=\res{(18/5;24/5)}$ m.
Le travail maximal est $\res{300\text{ J}}$.
\end{questions}
\label{figgeo-07-14-cor}
\figureGeo[corrige]{07-14}
\end{corrige}
\end{blocLivre}
\begin{blocLivre}{cor-7-15}
\begin{corrige}{15}{Une identité avec le milieu}
\demarche{Décomposer les vecteurs en passant par le milieu pour supprimer les termes croisés.}
\begin{questions}
\item Comme $I$ est le milieu, $\vect{IB}=-\vect{IA}$.
Les décompositions de Chasles donnent
\begin{align*}
\vect{MA}\cdot\vect{MB}
&=(\vect{MI}+\vect{IA})\cdot(\vect{MI}-\vect{IA})\\
&=MI^2-\vect{MI}\cdot\vect{IA}+\vect{IA}\cdot\vect{MI}-IA^2\\
&=MI^2-IA^2=\res{MI^2-AB^2/4}.
\end{align*}
La symétrie du produit scalaire fait disparaître les termes croisés.
\item Pour $A(0;0)$ et $B(6;0)$, $I=(3;0)$ et $AB^2/4=9$.
Le produit vaut 5 exactement lorsque $MI^2-9=5$, soit $MI=\sqrt{14}$.
Le lieu est donc le cercle de centre $\res{I(3;0)}$ et de rayon $\res{\sqrt{14}}$.
\item Comme $MI^2\geq0$, le produit est au moins $-9$.
Cette valeur est atteinte si et seulement si $MI=0$, donc uniquement en $\res{M=I}$.
\end{questions}
\label{figgeo-07-15-cor}
\figureGeo[corrige]{07-15}
\end{corrige}
\end{blocLivre}
